\section{Вывод}

В работе исследованы локальные и глобальные методы одномерной минимизации.
Эффективность методов оценивалась по длине итогового отрезка локализации,
числу вычислений функции и скорости уменьшения ошибки.

Основные результаты:

\begin{enumerate}
    \item \textbf{Метод Фибоначчи показал наилучшую эффективность при
    фиксированном числе вычислений.} Он оптимально использует доступный
    вычислительный бюджет, поскольку на каждой итерации сохраняет одно из
    ранее вычисленных значений функции.

    \item \textbf{Метод золотого сечения практически сопоставим с методом
    Фибоначчи.} Его преимущество заключается в более простой реализации и
    возможности остановки на любой итерации без предварительного выбора
    конечного числа шагов.

    \item \textbf{Метод дихотомии является надёжным, но менее экономичным.}
    На каждой итерации он требует двух новых вычислений функции, поэтому при
    одинаковом вычислительном бюджете обычно уступает методам Фибоначчи и
    золотого сечения.

    \item \textbf{Пассивный поиск является самым простым, но наименее
    эффективным локальным методом.} Для повышения точности ему приходится
    равномерно исследовать весь интервал, что приводит к большому числу
    вычислений функции.

    \item \textbf{Метод парабол эффективен для гладких функций.} Использование
    квадратичной аппроксимации позволяет быстро уточнять положение минимума.
    Однако метод чувствителен к выбору трёх точек и может работать нестабильно,
    если построенная парабола вырождается или её вершина выходит за пределы
    текущего интервала.

    \item \textbf{Локальные методы нельзя напрямую применять к
    мультимодальным функциям.} В этом случае они могут сойтись к одному из
    локальных минимумов, не обнаружив глобальный минимум.

    \item \textbf{Для мультимодальной функции метод ломаных является более
    содержательным, чем равномерный перебор.} Он распределяет вычисления
    адаптивно и использует оценку Липшица, поэтому способен уделять больше
    внимания наиболее перспективным участкам интервала.

    \item \textbf{Равномерный перебор отличается предсказуемостью, но требует
    большого числа вычислений.} Его преимущество состоит в простоте и
    равномерном исследовании всего интервала, однако он не использует
    информацию о форме функции.

    \item \textbf{Выбор метода должен зависеть от структуры целевой функции.}
    Для унимодальных функций наиболее рациональны методы Фибоначчи и золотого
    сечения, а для гладких функций с устойчивой аппроксимацией дополнительно
    эффективен метод парабол. Для мультимодальных функций следует применять
    методы глобального поиска.
\end{enumerate}

Таким образом, наиболее универсальным выбором среди рассмотренных локальных
методов является метод золотого сечения: он близок по эффективности к методу
Фибоначчи, но проще в реализации и допускает остановку в произвольный момент.
Метод Фибоначчи предпочтителен, когда заранее задано точное число вычислений.
Метод парабол способен дать лучший результат на гладких функциях, но требует
дополнительного контроля устойчивости. Для функций с несколькими экстремумами
необходимо использовать глобальные методы, поскольку локальная минимизация не
гарантирует нахождения глобального минимума.